\section{Introduction}\label{sec:intro}

Whether market pioneers are advantaged or burdened is an old question
that has been answered mostly on named markets and on the firms that
survived long enough to be remembered \citep{GolderTellis1993,
SemadeniAnderson2010}. The US trademark record offers a different sample.
Every firm that intends to sell a product under a mark files a dated
description of that product; the description is written by the millions
of firms that never patent as well as by those that do; and between the
fifth and sixth year of a registration the owner must swear that the
product is still in commerce or lose the mark. Whether a firm was early
can therefore be read from what every entrant wrote, and whether its
product lasted can be read from what every owner later proved. This paper
asks the pioneer's question on that record.

Being early is measured from the language itself. Each goods/services
description is reduced to themes by a topic model fitted once over all
13.99 million applications, and its theme mix is compared, by
Kullback--Leibler divergence, to what its own Nice class filed in the
five years before and the five years after it. The average of the two
divergences is the filing's \emph{atypicality}, how unusual it was for
its industry; their difference is its \emph{lead}, positive when the
industry's language subsequently moved toward the filing's. Lead is
computed from filings not yet made, so no one at the time could observe
it; it is a positional quantity, not a causal one. The construction, its
nomenclature, and the validation program that supports it --- three
partitions of the vocabulary, two reference weightings, five window
constructions, a second model seed, and the text duplication of a
boilerplate-heavy corpus --- are the subject of a companion paper, and
Section~\ref{sec:measure} summarizes what a reader of this one needs.

The answer is that the arrows are real. Among 3.10 million registrations
issued 2002--2018, marks whose language most led their industry were
cancelled for non-use at the five-year proof $1.4$ points more often than
marks whose language most lagged it, within Nice class and registration
year, with drafting and counsel held and errors clustered on the owner
($t = 8.3$); the raw contrast is $+2.3$ points, positive in 33 of 44
classes, and it holds under every partition of the vocabulary. Being
unusual carries no such cost: what protection atypical language appears
to carry belongs to filings that extend a theme established in another
class, and net of atypicality even those fail more.

The penalty is not constant, and where it moved is informative. It was
two to three times larger for marks filed in the late 1990s than for any
cohort since, reversed for marks filed 2000--2004, and returned from
2008, with the whole swing inside the four technology classes. Most of
the swing is between themes rather than within them: each era's leading
vocabulary was concentrated in one or two themes --- computer and
information services in the 1990s, downloadable software after 2008 ---
that failed at their own rate, while the penalty for being leading inside
a theme held at $+2.5$ to $+5$ points throughout. Web-based filings fail
hardest in each class's own first internet years and converge toward
their class within a decade. Joining a surging theme costs little on
average, and arriving early within a wave costs nothing further.

The record supplies two evaluators, and they read the text differently.
The examiner at registration rewards middling conformity --- completion
rises with atypicality and turns down only at the extreme --- and is
indifferent to lead, which did not yet exist to be read. The market at
the five-year proof is indifferent to conformity once theme importation
is held and penalizes having been early. Following the owners further,
no rung of financial success rewards it: the debut's language does not
predict a priced private round, SEC reporting and listing decline in
lead, and the firms holding most of the revenue the record can trace
debuted in language their industries already used. The paper also
records what applicants learn: those who refile an abandoned application
rewrite toward the class's typical description, from both tails, whoever
drafts the second attempt.

Section~\ref{sec:related} places the exercise against the pioneering,
optimal-distinctiveness and trademark-indicator literatures.
Section~\ref{sec:measure} states the measure and the record's steps.
Sections~\ref{sec:registration}--\ref{sec:success} follow the mark and
its owner through registration, the five-year proof, the waves that moved
the penalty, and the capital markets. Section~\ref{sec:discussion} reads
the stages together, and Section~\ref{sec:robust} summarizes what
changes under other choices.
